% Source: OpenAI2026AreaLaw, 09-amplification.tex, lines 289--295 and 493--575,
% eq:amplification-cut-clearance, eq:amplification-before-polar,
% eq:amplification-direct-error and eq:amplification-unitary-error, revision
% openai/math@adc7f1241b42e322a6451854ab7e4b4c146bf78a.

\section{The cut clearance and the local restoring unitary}

The proof of the amplification proposition ends with two deterministic
arguments. One keeps the enlarged collar inside the cut. The other turns a
good clock realization into a local unitary that nearly restores the
reference state. Neither uses the clocks or the spectral estimates.

\begin{theorem}[Enlarged collars stay inside the cut]
    \label{thm:area_law_amplification_cut_clearance}
    \lean{TNLean.PEPS.AreaLaw.exists_threshold_graphNeighborhood_collar_subset}
    \leanok
    \uses{def:area_law_radius_exponent}
    Let $\Lambda$ be a finite induced domain, let $A\subseteq\Lambda$, let
    $T\subset\mathbb Z^2$ be finite, and let $Z$ be the set of endpoints of the
    cut edges of $A$. Put $L=\lfloor n^{1-\varepsilon}\rfloor$ and
    $S_0=A\cap T_L$. For every $C$ there is $N$ such that for all $n\ge N$,
    every integer $r\le Cn^{(1-\varepsilon)/\alpha}$ and every $\Lambda$, $A$,
    $T$ with $\operatorname{dist}_\infty(T,Z)>L+n^\beta$ satisfy
    $N_{2r}^\Lambda(S_0)\subseteq A$. This is the last assertion of
    Proposition~10.2 of \cite{OpenAI2026AreaLaw}, given its radius bound.
\end{theorem}
\begin{proof}
    \leanok
    \uses{thm:area_law_uniform_collar_threshold,
        thm:area_law_graph_neighborhood_cut_clearance}
    Choose $N$ so that $L+2r<n^\beta$ for every admissible radius. Then
    $L+2r<L+n^\beta<\operatorname{dist}_\infty(T,Z)$, so a walk of length at
    most $2r$ from $S_0$ cannot cross the cut, and
    $N_{2r}^\Lambda(S_0)\subseteq A$.
\end{proof}

\begin{theorem}[Amplification from its radius and output bounds]
    \label{thm:area_law_amplification_of_radius_output}
    \lean{TNLean.PEPS.AreaLaw.amplificationBound_of_radius_output}
    \leanok
    \uses{def:area_law_amplification_bound}
    Fix $q,R,J,\Delta$. Suppose that, under the hypotheses of the amplification
    bound, for every $C_0>0$ and every integer $M\ge1$ there are $C_M$ and $N$
    such that at every $n\ge\max(2,N)$ some integer
    $r\le C_Mn^{(1-\varepsilon)/\alpha}$ satisfies $I_\Omega(X':J)\le C_Mn^{-M}$
    for every $X'\subseteq X$ and every
    $J\subseteq\Lambda\setminus N_{2r}^\Lambda(S_0)$. Then the amplification
    bound holds. This reduces Proposition~10.2 of \cite{OpenAI2026AreaLaw} to its
    radius and output bounds; it does not prove the proposition.
\end{theorem}
\begin{proof}
    \leanok
    \uses{thm:area_law_amplification_cut_clearance}
    Take the larger of the two thresholds. The cut-clearance clause follows from
    the preceding theorem applied to the radius supplied by the hypothesis.
\end{proof}

\begin{theorem}[Truncated products on a good realization]
    \label{thm:area_law_amplification_before_polar}
    \lean{TNLean.PEPS.AreaLaw.norm_truncated_restoring_le,
        TNLean.PEPS.AreaLaw.norm_truncated_adjoint_sub_le}
    \leanok
    Let $D$, $D_{\rm ret}$ and $L'$ be square matrices on one finite-dimensional
    Hilbert space, let $\zeta$ and $g$ be vectors, and let $M_0,\varepsilon$ be
    real. Suppose that $Dg=g$, $\|g\|\le1$, $\|\zeta\|\le M_0$,
    \begin{align*}
        \|D(\zeta-g)\| & \le\varepsilon,
        \qquad
        \|(D-D_{\rm ret})\zeta\|\le\varepsilon,
        \qquad
        M_0\|D_{\rm ret}-L'\|\le\varepsilon.
    \end{align*}
    Then $\|L'\zeta\|\le1+3\varepsilon$. If moreover $M_0\ge1$ and
    $D_{\rm ret}^\dagger w=w$ for a vector with $\|w\|\le1$, then
    $\|(L')^\dagger w-w\|\le\varepsilon$. With $g$ the ground component of the
    restoring vector, these are the two estimates that precede the polar
    decomposition in the proof of Proposition~10.2 of \cite{OpenAI2026AreaLaw}.
\end{theorem}
\begin{proof}
    \leanok
    Write
    $L'\zeta=g+D(\zeta-g)-(D-D_{\rm ret})\zeta-(D_{\rm ret}-L')\zeta$ and use
    $\|(D_{\rm ret}-L')\zeta\|\le\|D_{\rm ret}-L'\|M_0\le\varepsilon$. For the
    adjoint, $(L')^\dagger w-w=(L'-D_{\rm ret})^\dagger w$, and the adjoint has
    the same operator norm, at most $\varepsilon/M_0\le\varepsilon$.
\end{proof}

\begin{theorem}[From an accurate adjoint to a local unitary]
    \label{thm:area_law_amplification_unitary_error}
    \lean{TNLean.PEPS.AreaLaw.norm_sub_sq_le_of_adjoint_error,
        TNLean.PEPS.AreaLaw.exists_unitary_kronecker_sub_le}
    \leanok
    Let $R'$ be a square matrix on a finite-dimensional Hilbert space
    $\mathcal H$, and let $\phi,\eta$ be unit vectors in
    $\mathcal H\otimes\mathcal K$ for a finite-dimensional space $\mathcal K$.
    If $\|(R'\otimes I)^\dagger\eta-\phi\|\le a$ and
    $\|(R'\otimes I)\phi\|\le1+b$, then
    \begin{align*}
        \|(R'\otimes I)\phi-\eta\|^2 & \le2a+2b+b^2,
    \end{align*}
    and there is a unitary $U$ on $\mathcal H$ with
    $\|(U\otimes I)\phi-\eta\|\le a+\sqrt{2a+2b+b^2}$. No invertibility of
    $R'$ is assumed. These are the direct and unitary errors in the proof of
    Proposition~10.2 of \cite{OpenAI2026AreaLaw}.
\end{theorem}
\begin{proof}
    \leanok
    The real part of $\langle\eta,(R'\otimes I)\phi\rangle$ equals that of
    $\langle(R'\otimes I)^\dagger\eta,\phi\rangle$, which is at least $1-a$.
    Expanding the square gives the first bound. Take $U$ to be a unitary polar
    factor of $R'$, so that $R'=UH_+$ with $H_+\ge0$. Since $I+H_+\otimes I$
    does not decrease norms, the distance from $(U\otimes I)\phi$ to $\eta$ is at
    most the sum of the direct and adjoint errors.
\end{proof}

\begin{theorem}[A local unitary from a good realization]
    \label{thm:area_law_amplification_good_realization}
    \lean{TNLean.PEPS.AreaLaw.exists_unitary_kronecker_sub_le_of_good_realization}
    \leanok
    Let $R_0$ and $L'$ be square matrices on a finite-dimensional space
    $\mathcal H$, and let $D$, $D_{\rm ret}$ act on $\mathcal H\otimes\mathcal K$.
    Let $\phi,\eta$ be unit vectors, put $\zeta=(R_0\otimes I)\phi$, and suppose
    that $\|\zeta\|\le M_0$, that $D$ fixes a vector $g$ with $\|g\|\le1$, that
    \begin{align*}
        \|D(\zeta-g)\| & \le\varepsilon,
        \qquad
        \|(D-D_{\rm ret})\zeta\|\le\varepsilon,
        \qquad
        M_0\|D_{\rm ret}-L'\otimes I\|\le\varepsilon,
    \end{align*}
    and that $\|((L'R_0)\otimes I)^\dagger\eta-\phi\|\le a$. Then some unitary
    $U$ on $\mathcal H$ satisfies
    $\|(U\otimes I)\phi-\eta\|\le a+\sqrt{2a+6\varepsilon+9\varepsilon^2}$.
\end{theorem}
\begin{proof}
    \leanok
    \uses{thm:area_law_amplification_before_polar,
        thm:area_law_amplification_unitary_error}
    The first truncation estimate gives $\|(L'R_0\otimes I)\phi\|\le1+3\varepsilon$.
    Apply the preceding theorem with $b=3\varepsilon$.
\end{proof}
